\subsection{循环扩张}

定义 3. --- 若 K 的扩张 E 是伽罗瓦扩张且其伽罗瓦群是循环群，则称 E 为循环扩张。

例子. --- 1) 每个素数次的伽罗瓦扩张都是循环的，因为每个素数阶p的有限群G都是循环的（对于G中的每个元素 \(x \neq 1\) ，其阶为p，因此生成G）。

\begin{enumerate}
\def\labelenumi{\arabic{enumi})}
\setcounter{enumi}{1}
\item
  设 \(\mathrm{F}(\mathbf{X}) = \mathbf{X}^2 + a\mathbf{X} + b\) 是 \(\mathbf{K}[\mathbf{X}]\) 中的一个不可约多项式。 \(\mathrm{F}(\mathbf{X})\) 不可分的唯一情形是 \(\mathbf{K}\) 的特征为2且 \(a = 0\) 的情形。我们将这种情形搁置一旁；设 \(\mathbf{E}\) 是 \(\mathbf{K}\) 的、由一个根 \(x\) （ \(\mathrm{F}(\mathbf{X})\) 的根）生成的扩张。我们有 \([\mathrm{E}:\mathrm{K}] = 2\) 和 \(\mathrm{F}(\mathbf{X}) = (\mathbf{X} - x)(\mathbf{X} + a + x)\) ，因此 \(\mathbf{E}\) 是 \(\mathbf{K}\) 的伽罗瓦扩张。 \(\mathbf{E}\) 在 \(\mathbf{K}\) 上的伽罗瓦群是2阶的，从而是循环的。
\item
  设 F 是一个域，且 \(\sigma\) 是一个有限阶自同构。 \(\sigma\) 的不变量域 E 也是由 \(\sigma\) 生成的 n 阶循环群的不变域，因此（V，第66页，定理3）F 是 E 的一个 n 次循环扩张。
\end{enumerate}

我们知道（I，第50页）循环群的每个子群和每个商群都是循环群。因此（V，第68页，推论4），若E是域K的n次循环扩张，则E的每个子扩张F在K上是循环的，且E在F上是循环的。对于n的每个因子d，存在唯一一个包含于E中的、在K上次数为d的子域F：因为在n阶循环群中存在唯一一个指数为d的子群。

定理 3（希尔伯特）。--- 设 E 是 K 的一个循环扩张，且令 \(\sigma\) 为 E 在 K 上的伽罗瓦群 \(\Gamma\) 的一个生成元。

\begin{enumerate}
\def\labelenumi{\alph{enumi})}
\item
  为使元素 \(x \in \mathbb{E}\) 满足 \(\mathrm{N}_{\mathrm{E} / \mathrm{K}}(x) = 1\) ，必须且只需存在 \(y \in \mathrm{E}^*\) 使得 \(x = y / \sigma(y)\) ；于是每个 \(y_1 \in \mathrm{E}^*\) （满足 \(x = y_1 / \sigma(y_1)\) ）都具有形式 \(\lambda y\) ，其中 \(\lambda \in \mathbf{K}^*\) 。
\item
  为使元素 \(x \in E\) 满足 \(\operatorname{Tr}_{\operatorname{E}/\operatorname{K}}(x) = 0\) ，必须且只需存在 \(z \in E\) 使得 \(x = z - \sigma(z)\) ；于是每个 \(z_{1} \in E\) （满足 \(x = z_{1} - \sigma(z_{1})\) ）都具有形式 \(z + \mu\) ，其中 \(\mu \in K\) 。
\end{enumerate}

让我们首先证明一个引理。

引理4. --- 设 \(\Gamma\) 是一个阶为 \(n\) 的循环群， \(\sigma\) 是 \(\Gamma\) 的一个生成元，且 M 是一个交换群， \(\Gamma\) 按照规则 \(\gamma\) . \((m + m') = \gamma\) . \(m + \gamma\) . \(m'\) 作用在 M 上（对所有 \(\gamma \in \Gamma\) 和 \(m\) , \(m' \in M\) ）。设 Z 为所有从 \(\Gamma\) 到 M 中且满足下述关系的映射

\[
f \left(\tau \tau^ {\prime}\right) = f (\tau) + \tau . f \left(\tau^ {\prime}\right) \quad \text { for } \quad \tau , \tau^ {\prime} \text { in } \Gamma .\tag{10}
\]

令 \(u(f) = f(\sigma)\) （对 \(f \in \mathbf{Z}\) ）和 \(t(m) = \sum_{\tau \in \Gamma} \tau \cdot m\) （对 \(m \in \mathbf{M}\) ）。则序列

\[
0 \to \mathbf {Z} \xrightarrow {u} \mathbf {M} \xrightarrow {t} \mathbf {M}
\]

是正合的。

设 \(f \in Z\) ；在（10）中取 \(\tau = \tau' = 1\) ，我们得到 \(f(1) = 0\) 。进一步地，通过对 \(m \geqslant 0\) 作归纳，我们推出关系式

\[
f (\sigma^ {m}) = f (\sigma) + \sigma \cdot f (\sigma) + \dots + \sigma^ {m - 2} \cdot f (\sigma) + \sigma^ {m - 1} \cdot f (\sigma).\tag{11}
\]

我们有 \(\sigma^n = 1\) ，从而 \(f(\sigma^n) = 0\) ；前述关系式取 \(m = n\) 时等价于等式 \(t(u(f)) = 0\) ，从而 \(\operatorname{Im} u \subset \ker t\) 。此外，由(11)可得 \(\operatorname{Ker} u = 0\) 。

设 \(m \in M\) 使得 \(t(m) = 0\) ，即 \(m + \sigma \cdot m + \cdots + \sigma^{n-1} \cdot m = 0\) 。我们定义映射 \(f\) （从 \(\Gamma\) 到 \(M\) 中）如下

\[
f \left(\sigma^ {j}\right) = m + \sigma . m + \dots + \sigma^ {j - 2}. m + \sigma^ {j - 1}. m\tag{12}
\]

（对 \(0 \leqslant j \leqslant n - 1\) ）。把关系式（10）的验证留给读者。我们显然有 \(m = f(\sigma)\) ，从而 \(\operatorname{Im} u \supset \operatorname{Ker} t\) 。

既然引理已经证明，取 \(y \in E^{*}\) 且 \(x = y / \sigma(y)\) ；我们有 \(\mathrm{N}_{\mathrm{E}/\mathrm{K}}(\sigma(y)) = \mathrm{N}_{\mathrm{E}/\mathrm{K}}(y)\) ，从而 \(\mathrm{N}_{\mathrm{E}/\mathrm{K}}(x) = 1\) 。反之，设 \(E^{*}\) 中的 x 使得 \(\mathrm{N}_{\mathrm{E}/\mathrm{K}}(x) = 1\) ；由引理4（应用于 \(M = E^{*}\) ），存在一族元素 \((c_{\tau})_{\tau \in \Gamma}\) （属于 \(E^{*}\) ），满足关系 \(c_{\tau\tau'} = c_{\tau} \cdot \tau(c_{\tau'})\) （其中 \(\tau\) , \(\tau'\) 属于 \(\Gamma\) ）以及 \(c_{\sigma} = x\) 。由命题9（V，第65页）的推论1，存在 \(y \in E^{*}\) 使得 \(c_{\tau} = y / \tau(y)\) 对所有 \(\tau \in \Gamma\) 成立，从而特别地 \(x = c_{\sigma} = y / \sigma(y)\) 。若 \(y_{1} \in E^{*}\) 满足 \(x = y_{1} / \sigma(y_{1})\) ，则我们有

\[
\sigma \left(y _ {1} y ^ {- 1}\right) = y _ {1} y ^ {- 1},
\]

因此 \(y_{1}y^{-1}\) 属于 \(K^{*}\) ，因为 \(\sigma\) 生成 E 在 K 上的伽罗瓦群。这证明了 a）。

断言b）以类似方式由命题9（V，第65页）的推论2得出。

